\section{Introduction}
\label{sec:introduction}

Multi-agent AI systems that recursively delegate objectives face a structural correctness failure we term the \textbf{autonomous drift problem}. In recursive spawning architectures, an agent may initialize one or more child agents to accomplish a delegated objective. These child agents may, in turn, initialize further descendants. As delegation chains grow in depth and branching factor, they lose traceable connectivity to the original authorizing human principal. The result is an agent that continues executing — actively spawning descendants, exercising delegated authority — while being operationally disconnected from any human accountability chain.

We formalize this failure state as the \textbf{drift triad}:

\begin{equation}
\text{DriftTriad}(n) \;=\; \bigl\langle \text{orphaned}(n),\; \text{drifted}(n),\; \text{operating}(n) \bigr\rangle .
\label{eq:intro-drift-triad}
\end{equation}

An agent is \textbf{orphaned} when its parent has terminated or become unreachable. An agent is \textbf{drifted} when its delegation chain does not reach any human principal at its origin — no principal stands behind the authority it exercises. An agent \textbf{operates} when it continues executing its event loop. The conjunction — orphaned \emph{and} unanchored \emph{and} actively running — is the specific failure state this paper addresses. Critically, drift is heritable: a drifted agent that spawns a child propagates the condition to all descendants without detection, producing what we call a \textbf{drift cascade}. A single undetected drift event at depth $d$ with branching factor $b$ yields up to $b^d$ drifted descendants. This is not a governance concern solvable by policy conventions. It is a structural correctness failure.

The field has attempted to address drift through two primary mechanisms. \textbf{Time-to-live} (TTL) constraints limit the operational lifetime of an agent or the number of delegation hops before workflow termination. \textbf{Depth limits} cap the maximum nesting depth of recursive spawning, restricting how many generations can exist simultaneously. Both impose an external, temporal bound on system behavior. Neither addresses the fundamental structural problem: in existing systems, when a sub-agent spawns a further sub-agent, the parent-child relationship is stored as a mutable reference. It can be modified, overwritten, or lost. TTL expiration does not preserve lineage — it terminates. Depth limits do not encode authorization scope propagation — they truncate. As a result, existing systems can enforce \emph{when} a system stops, but not \emph{what} it did or \emph{who authorized it}.

The deeper insufficiency of depth limits is that they constrain chain \emph{length} — a topological property of the spawn graph — while drift is fundamentally about authorization \emph{scope}: whether each node's operating envelope is a legitimate descendant refinement of a human principal's intent. A drifted agent at depth $1$ with scope expansion authority can produce consequences that exceed the human anchor's intent as severely as a drifted agent at depth $D_{\max}-1$. No choice of $D_{\max}$ prevents scope creep drift without an additional scope refinement constraint. We prove in Section~\ref{sec:autonomous-drift} that drift prevention requires the conjunction of two independent constraints:

\begin{equation}
\text{Drift Prevention} \;\iff\; \bigl( \text{depth}(C) \leq D_{\max} \bigr) \;\land\; \bigl( \forall i:\; R(n_{i+1}) \subseteq R(n_i) \bigr).
\label{eq:intro-drift-prevention}
\end{equation}

Depth limits enforce only the first conjunct. The second — scope monotonicity — is structurally absent from TTL- and depth-limit-based systems.

The core insight of this paper is that the autonomous drift problem is structurally soluble: it is not an inherent consequence of recursive spawning, but a consequence of \emph{mutable} parent provenance. If every agent's parent reference is fixed immutably at spawn time and embedded in the agent's foundational identity — not stored as a mutable attribute — then the agent cannot voluntarily or inadvertently dissociate from its chain of authorization. We formalize this through the \textbf{immutable parent pointer chain}: a cryptographically signed, append-only record of ancestor references sealed at agent creation. Any attempt to modify an ancestor record after the fact is rejected by the protocol. The \textbf{BX3 Framework} provides the architectural substrate that makes this construction enforceable in operational runtime.

BX3 organizes agent governance into three interdependent layers. The \textbf{Purpose Layer} carries the terminal accountability anchor for every agent: each root agent is initialized with a declared \textit{purpose clause} defining the boundaries of its operational scope, and no child agent may operate outside its parent's purpose semantic envelope. The \textbf{Bounds Engine} enforces structural constraints — maximum depth $D_{\max}$ and resource budgets — preventing unbounded recursion as a practical and accountability matter. The \textbf{Fact Layer} is the append-only, tamper-evident forensic ledger that records every spawn event, delegation action, purpose clause binding, and agent state transition. Its pre-commit hook is the enforcement point: any spawn event that violates the depth bound or the scope refinement constraint is rejected before the child agent enters operational state. No layer alone is sufficient. Together, they make the drift triad — orphaned, drifted, and operating simultaneously — architecturally impossible, not merely statistically unlikely.

The remainder of this paper is organized as follows. Section~\ref{sec:autonomous-drift} formalizes the drift problem, establishes the drift triad and its constituent sub-modes, identifies three structural enabling conditions, catalogs four concrete failure modes, and proves that depth limits alone are insufficient. Section~\ref{sec:bx3-integration} presents the integration of Recursive Spawning with the BX3 Purpose Layer, Bounds Engine, and Fact Layer, showing how they collectively enforce drift prevention as a structural invariant. Section~\ref{sec:rs-protocol} specifies the Recursive Spawning Protocol (RSP) — the mandatory five-stage spawn sequence executed across all three BX3 layers at every spawn event — and proves three correctness properties: drift prevention, chain integrity, and finite termination. Section~\ref{sec:experimental} presents experimental evaluation measuring drift resistance under adversarial propagation scenarios. Sections~\ref{sec:limitations} and~\ref{sec:conclusion} discuss limitations, ethical implications, and open problems.